% !TeX root = ../thesis.tex
\chapter{Appendices to Chapter~\ref{ch:netcomi}}
\label{sec:app:NetCoMi}

\begin{figure}[H]
  \centering
  \includegraphics[width=0.85\textwidth]{figures/app_NetCoMi/netcomi_network_gcms-1.pdf}
  \caption{\textbf{Graphlet Correlation Matrices for illustrative network models.} Each row shows an unweighted network model with $117$ nodes and its corresponding Graphlet Correlation Matrix (GCM), computed from the $11$ non-redundant graphlet orbits used in NetCoMi. Across the three standard models, the GCMs become more structured as the networks move away from a homogeneous random baseline. The random Erd\H{o}s--R\'enyi network mainly shows degree-driven positive orbit correlations, the Barab\'asi--Albert network highlights hub-related dependencies, and the stochastic block model emphasizes module-related dependencies. The figure illustrates that networks with different generative mechanisms can yield characteristic graphlet-correlation patterns, supporting the use of GCM heatmaps as comparative summaries of local network topology.}
  \label{fig:app:5net:gcm_network_models}
\end{figure}
